\setlength{\absparsep}{18pt}
\begin{resumo}
Software embarcado orientado a interrupções está sujeito a \textit{data races} e
violações de atomicidade: uma rotina de tratamento de interrupção pode
interromper o código principal entre dois acessos a uma variável compartilhada
e corromper seu valor. As ferramentas da literatura priorizam a precisão, medem
mal o recall e tratam essas duas classes de defeito separadamente; apenas a mais
antiga propõe correções. Em paralelo, trabalhos recentes usam modelos de
linguagem para triar e reparar alertas de analisadores estáticos, mas todos
tratam defeitos sequenciais. Este trabalho propõe uma ferramenta que ocupa essa
interseção. Uma análise estática sobre a representação intermediária do LLVM
encontra os candidatos das duas classes sem falsos negativos, em relação a uma
lista explícita de premissas; um solver SMT descarta apenas o que prova
impossível; e um modelo de linguagem ordena, explica e propõe correções para os
candidatos restantes, sem nunca descartá-los. A avaliação usa o benchmark
Racebench, com gabarito recontado e corrigido (48 defeitos e 38 armadilhas),
métricas que não podem ser melhoradas descartando candidatos (portão de recall,
rejeição de armadilhas e \textit{Inspection Ratio}) e uma ablação das técnicas de
\textit{prompt}. Esta monografia parcial apresenta a fundamentação, a revisão da
literatura, a proposta, a metodologia de avaliação e o cronograma do trabalho.

Palavras-chave: concorrência. interrupções. software embarcado.
análise estática. modelos de linguagem.
\end{resumo}
